% Generated draft: requires mathematical and source review.
\documentclass[11pt]{article}
\usepackage[margin=25mm]{geometry}
\usepackage{fontspec}
\setmainfont{Latin Modern Roman}
\usepackage{amsmath,amssymb,mathtools,unicode-math}
\setmathfont{Latin Modern Math}
\usepackage{enumitem,needspace,xurl,hyperref}
\setlength{\parindent}{0pt}
\setlength{\parskip}{5pt}
\setlength{\emergencystretch}{3em}
\newcommand{\N}{\mathbb{N}}
\newcommand{\Z}{\mathbb{Z}}
\newcommand{\Q}{\mathbb{Q}}
\newcommand{\R}{\mathbb{R}}
\newcommand{\C}{\mathbb{C}}
\newcommand{\F}{\mathbb{F}}
\newcommand{\A}{\mathbb{A}}
\newcommand{\PP}{\mathbb{P}}
\newcommand{\E}{\mathbb{E}}
\newcommand{\eps}{\varepsilon}
\newcommand{\dd}{\,\mathrm{d}}
\begin{document}
\section*{UnsolvedMath ID: 30003836}
\textbf{Source problem code:} catalogue entry
\subsection{Definitions and mathematical statement}
The supplied record names the Monod groups \(H(A)\) and the Lodha--Moore group \(G_0\), but provides neither their definitions nor the admissible parameters \(A\). Consequently, the following describes the mathematical targets conditionally on identifying those groups and any intended actions; it is not a complete specification of a single problem.

For \(X=\R\) or \(X=[0,1]\), let \(\operatorname{Homeo}(X)\) denote the group of homeomorphisms of \(X\), and let \(\operatorname{Diff}^{1}(X)\) denote the group of bijections of \(X\) whose maps and inverses are continuously differentiable. For the interval, use the convention that derivatives extend continuously to the endpoints as one-sided derivatives. The source does not specify its endpoint convention or whether only orientation-preserving maps are intended; these choices require confirmation. The interval \([0,1]\) represents a nondegenerate compact interval by affine change of coordinates, assuming this is what the source means by a closed interval.

A faithful \(C^1\) action of a group \(K\) on \(X\) is an injective homomorphism
\[
\rho:K\longrightarrow\operatorname{Diff}^{1}(X).
\]
For a specified action \(\alpha:K\to\operatorname{Homeo}(X)\), say that \(\alpha\) is \(C^1\)-smoothable if there exists \(h\in\operatorname{Homeo}(X)\) such that
\[
h\alpha(g)h^{-1}\in\operatorname{Diff}^{1}(X)
\qquad\text{for every }g\in K.
\]
This definition concerns a specified action, whereas the existence of a faithful \(C^1\) action concerns the abstract group. The supplied record does not identify the actions to be smoothed and does not justify identifying these two questions.

The recorded questions are: for the intended Monod groups \(H(A)\) and the intended Lodha--Moore group \(G_0\), are their intended actions \(C^1\)-smoothable, and do these groups admit faithful \(C^1\) actions on \(\R\) or on a nondegenerate closed interval? For each intended group \(K\), the latter questions take the forms
\[
\exists\,\rho_{\R}:K\hookrightarrow\operatorname{Diff}^{1}(\R),
\qquad
\exists\,\rho_I:K\hookrightarrow\operatorname{Diff}^{1}([0,1]).
\]
They are displayed separately because the source does not specify whether its word ``or'' requests either realization, both realization questions, or a comparison with smoothability. The domain and quantifier for \(A\), the group definitions, and the designated actions remain unresolved.
\subsection{Short English statement}
The record asks about \(C^1\) smoothability and faithful differentiable actions of Monod groups \(H(A)\) and the Lodha--Moore group \(G_0\) on the real line or a closed interval. A precise statement requires the missing group definitions, the allowed parameters \(A\), and the actions intended for smoothing. Smoothability of a specified action and existence of some faithful action are distinct targets. The supplied open status is dated curation, not an independently verified conclusion.
\subsection{Sources}
\begin{itemize}
\item Local UnsolvedMath registry, \texttt{data/unsolvedmath/problems.json}, permanent ID \(30003836\), source problem code \texttt{catalogue entry}, title: \(C^1\) Smoothability of Monod and Lodha--Moore Groups. The supplied collection is Oberwolfach Reports Open Problems. The record gives proposed year 2018 and names no individual proposer.
\item Supplied source citation: \emph{Cohomological and Metric Properties of Groups of Homeomorphisms of R} (2019), Oberwolfach Reports, pp.\ 1579--1633, DOI \texttt{10.4171/owr/2018/26}. \url{https://doi.org/10.4171/owr/2018/26}. The background describes Oberwolfach Reports as documenting mathematical workshops and their open-problem discussions.
\item Supplied literature source, labelled ``Oberwolfach Report 2019'': \url{https://ems.press/content/serial-article-files/46748?nt=1}.
\item Supplied literature source, attributed to Lodha--Matte Bon--Triestino: \url{https://doi.org/10.1112/topo.12151}. No article title or specific supporting result is supplied.
\item Literature assessment dated exactly \texttt{2026-08-21}: the supplied curation labels the problem open, states that the groups are not identified in the truncated record, and says available smoothability results cover special classes only, leaving the exact question unresolved. This is supplied dated curation, not independently verified literature assessment. Its instruction to verify time-sensitive status and cited claims independently remains applicable.
\item The supplied statement status is \texttt{corrected\_verified}; its verification note claims that manual semantic, typographic, and LaTeX review found the clean statement self-contained and faithful to the target. That source-provided claim does not resolve the missing definitions and qualifications identified here. No linked source has been independently consulted for this draft.
\end{itemize}
\end{document}
